\hypertarget{rule-selection}{%
\section{Rule selection}\label{method-standards-methods-rule-selection}}

The proposed selection procedure distinguishes operative authority from historical source support. The retained sources include official AASB compilations, the original EU incorporating regulation and official ASBE lease and income-tax texts \citep{aasb9dated,aasb15dated,aasb16dated,aasb112dated,eu20231803,mofasbe21,moflease2018,mofasbe18,mofinterpretation16}. Historical US lease additions are separately identified in the retained amendment source \citep{fasbasu201602}.

AuthorityRequirement contains issuer/framework, jurisdiction, entity class, annual-period/application interval, paragraph and local source Ref, incorporating instrument, amendment chain, transition, elections and unresolved dependencies. RuleSelection binds those requirements to the accounting event and evidence view. The algorithm evaluates chronological applicability first, then scope and facts, then available elections/transition, then measurement/presentation. If two supported rules can apply, retain the explicit alternative and required election; if one is contradicted by facts, mark ineligible. Missing current authority gives unresolved, never eligibility through a familiar rule name.

An unchanged historical report remains an assertion under its own basis and vintage. A harmonized comparison creates an adjustment book and reconciliation, with changed recognition/measurement explained. Label mapping alone cannot convert US GAAP to IFRS or ASBE. The four clocks and source version remain part of each rule dependency; later knowledge cannot leak into an original as-of view. Transition entries reconcile opening assets, liabilities, tax and equity, rather than changing old source bytes or plugging an unexplained amount.

\hypertarget{ifrs-and-incorporating-jurisdictions}{%
\section{IFRS and incorporating jurisdictions}\label{method-standards-methods-ifrs-and-incorporating-jurisdictions}}

The retained official AASB texts provide explicit IAS/IFRS incorporation for their stated entity/version scope. AASB 9/15 compilation dates are 31 December 2022 with the source covers' application window; AASB 16 compilation 8 is 31 December 2023, AASB 112 compilation 7 is 29 September 2023. Their filenames are not date evidence. AASB 102's older retained edition is selected historical support for inventory, not current 2024 applicability certification. The retained EU 2023/1803 original provides official EU-adopted IFRS 3/10/IAS 19/21/36/37/IFRS 2/13 passages inspected in this method investigation; its historical adoption scope and later amendments remain separate. The retained amendment investigation supplies additional source evidence, not a blanket acceptance of all rules in these standards.

Initial real-method admission is therefore claim-specific: the selected revenue, lease/tax, financing, credit, group/FX and insurance cases must resolve each named applicability prerequisite. For synthetic demonstrations, declare the chosen dated provisions and assumptions and use the synthetic book mode. No assertion of whole-IFRS compliance follows from those entries.

The base lease in the saved dossiers (two annual 121 payments, 10\% rate, no qualifying exemptions or extra costs) has liability 210 initially, 110 after year one, 0 after year two. IFRS-equivalent cost-model ROU is 210/105/0; expenses 126 then 116. Income tax needs its own legal deductions and gross temporary differences, as the connected corporate case demonstrates. Revenue timing, inventory cost formula, financial instrument classification, control and contract boundaries are required rule predicates, not accounting options chosen after seeing which outcome fits best.

\hypertarget{us-gaap-proposed-methods-and-authority-limitation}{%
\section{US GAAP proposed methods and authority limitation}\label{method-standards-methods-us-gaap-proposed-methods-and-authority-limitation}}

The saved official FASB ASU 2016-02 Section A and later 2020-05/2021-09 amendments contain inspected historical ASC 842 additions. Their covers explicitly say an ASU is not authoritative. Current Codification access returned 403 again in this source investigation; the proposed methods cannot be marked current-authority-admitted. ASU amendments, historical FAS 109, issuer accounting policy and a model's memory are all insufficient substitutes for the operative ASC text and full applicable amendment chain. 

The same lease facts include ten-year remaining equipment life, fair value 1000, alternative use to lessor, no ownership transfer, no purchase option, two-year term and PV 210. Under the inspected historical 842-10-25-2/25-3 provisions, the stated facts support the proposed operating classification. Historical 842-20-25-6 and 35-3 yield single lease expense 121 per year, liability 210/110/0 and ROU 210/ 110/0. First year debit lease expense 121/credit liability 21/credit ROU 100; payment debits liability 121/credits cash 121. Next expense 121 decomposes into liability interest 11 and ROU reduction 110. The finance comparison uses the separate depreciation 105 plus interest 21/11 path. Changing classification facts is not a free election. The 75\%/90\% heuristics in the historical guidance cannot become universal requirements. This complete numerical contrast is useful despite the remaining current-authority limitation.

Proposed ASC 740 temporary-difference treatment retains tax base, reversal dates, enacted rates, valuation allowance and source-specific allocation/presentation. The saved historical FAS 109 and ASU 2019-12 provide inspected enactment support, not a completed current ASC 740 implementation. Under the matched PPE case, book 800/tax 600 at assumed 25\% yields DTL 50 and current tax 75 on taxable 300. A future 30\% rate that is merely substantively enacted but not legally enacted cannot be imported from the AASB 112 rule into this proposed US treatment; the declared 2025 closing temporary difference is 300, giving DTL 75 at the enacted 25\% rate versus DTL 90 at 30\%. This deferred-tax comparison is separate from current tax 75 on taxable profit 300. Actual legal enactment remains a separate jurisdictional fact.

The proposed CECL method for in-scope amortised-cost assets estimates lifetime expected cash shortfalls from origination using past/current/forecast evidence, with a documented reversion method beyond a reasonable/supportable forecast horizon and applicable instrument exceptions. It retains allowance roll-forwards, write-offs, recoveries, purchased-credit-deteriorated handling and off-balance commitments separately. It must not copy IFRS 9's 12-month/SICR staging to the same record. A candidate contrast: lifetime discounted loss 6 with 2 attributable to default events in the next year gives IFRS general Stage 1 candidate 2 versus CECL lifetime candidate 6, conditional on scope and method support. No current ASC 326 text was retrieved, and the attempted official amendment path returned 403; this contrast remains an explicitly unadmitted proposed reading, not a source-verified US requirement or posting example. Full current ASC 326, insurance 944, consolidation 810, FX 830 and relevant later amendments are specific internal source tasks before corresponding real admission.

\hypertarget{asbe-official-lease-support-and-remaining-tax-sources}{%
\section{ASBE lease adoption and deferred tax}\label{method-standards-methods-asbe-official-lease-support-and-remaining-tax-sources}}

For an ASBE lease, selecting the lease standard and selecting the related income-tax treatment require different application dates. Article 68 of the revised ASBE 21 states a general date of 1 January 2019, but the Ministry of Finance's 2018 issuance notice specifies the enterprise groups. Paragraph I requires application from that date for enterprises listed both domestically and overseas, and for overseas-listed enterprises preparing financial statements under IFRS or ASBE. Other enterprises applying ASBE must apply the revision from 1 January 2021 \citep[issuance notice, para.\ I]{moflease2018}.

Paragraph II permits earlier application for enterprises whose parent or subsidiary is listed overseas and prepares its overseas financial statements under IFRS or ASBE. Their early date cannot precede simultaneous application of the specified 2017 financial-instruments and revenue standards. Paragraph III discontinues the 2006 lease standard and its application guidance when the enterprise applies the revision. Thus a 2020 reporting date alone does not establish use of the revised lease standard for an enterprise in the second group: its qualifying early adoption would need evidence. The retained official Chinese texts include the 2018 notice and attachment and the Ministry's later ASBE 21 page, published on 28 October 2019 \citep{moflease2018,mofasbe21}. They replace the earlier secondary reproduction as support for the inspected wording.

Articles 4--8 require control of identified use, economic benefits and direction; article 6 includes both conditions for substantive substitution and the express presumption when difficult to determine. Articles 14--18 specify commencement, ROU cost and PV, using implicit rate or incremental borrowing rate if the former cannot be determined. Articles 20--23 separate cost-model depreciation and periodic interest; articles 25--29 distinguish reassessment/index changes/modifications and their discount-rate rules. Articles 30--33 define the exemptions; a two-year term is not short-term. Articles 60--67 specify transition choices and disclosures; the chosen transition cannot be inferred from a present closing liability \citep[attached standard, arts.\ 4--8, 14--33, 60--67]{moflease2018}.

For the declared base lease, these inspected recognition/measurement provisions support proposed ROU 210/105/0 and liability 210/110/0 with expense 126/116, under the same stated facts, depreciation assumption and fixed rule version. A same-asset modification is tested against 28--29, rather than labeled a separate lease without increased scope and commensurate consideration. A 2019 first-adopter using article 61(2) records cumulative opening changes without adjusting comparatives, while 61(1) is retrospective; comparison datasets must retain that distinction and the entity's actual election.

The tax question is why equal opening lease balances can still require separate deferred-tax assets and liabilities. The Ministry's original ASBE 18 defines tax bases and temporary differences in articles 4--9 and contains initial-recognition exceptions in articles 11 and 13. Article 13 also limits a deferred-tax asset to the probable taxable profit available to use the deduction. Articles 16--23 govern measurement, recovery or settlement assumptions, subsequent review, allocation and noncurrent classification \citep{mofasbe18}. The original text must be read with later interpretations.

Interpretation 16 removes the initial-recognition exemption for a single transaction that is not a business combination, initially affects neither accounting profit nor taxable income or loss, and creates equal taxable and deductible temporary differences. Section I(a) expressly includes qualifying initial lease liabilities and right-of-use assets; the corresponding deferred-tax liabilities and assets are recognized separately under ASBE 18. Equal carrying amounts alone do not establish equal temporary differences, and the asset's recoverability condition still applies. Section IV makes this treatment effective from 1 January 2023; Section I(b) permits early application from the year of issuance, 2022, with disclosure \citep[secs.\ I(a)--(b), IV]{mofinterpretation16}. This tax date is distinct from the enterprise's lease-standard adoption date. The operative text now supports the treatment previously illustrated by the retained official 2023 tax application case.

For the synthetic two-year lease, retain the assumed payments of 121, 10\% discount rate and straight-line depreciation. All amounts use the same illustrative monetary unit. Assume that tax deductions are allocated to the liability, both tax bases are zero, the applicable tax rate is 25\%, and sufficient taxable profit satisfies the deferred-tax asset recognition condition. Also assume that the transaction meets Section I(a)'s conditions and falls within its application period. At commencement the ROU asset and liability are each 210. The taxable ROU difference gives a deferred-tax liability (DTL) of 52.5; the deductible liability difference gives a deferred-tax asset (DTA) of 52.5. The entry debits DTA and credits DTL by those amounts. At the first annual close, liability 110 and ROU 105 give DTA 27.5 and DTL 26.25. Reducing the DTA by 25 creates deferred-tax expense; reducing the DTL by 26.25 creates a benefit, leaving net benefit 1.25. Both balances reverse by expiry. Applying the old exemption to omit both opening deferred balances is wrong for these qualifying facts and dates: it conceals the required gross balances and the later tax effect.

Transition also needs more than an opening journal entry. Section I(b) requires adjustments for qualifying transactions between the earliest comparative opening and application. For qualifying existing balances at that opening, cumulative effects adjust opening retained earnings and the relevant financial-statement items. Consider a separate transition example with cash 500, ROU 105, lease liability 110 and equity 495, with no deferred tax yet recognized. Under the same assumed tax bases, rate and recoverability, recognition debits DTA 27.5, credits DTL 26.25 and credits opening retained earnings 1.25. The resulting balance sheet is
\[
\underbrace{500+105+27.5}_{\text{assets }632.5}
=\underbrace{110+26.25}_{\text{liabilities }136.25}
+\underbrace{495+1.25}_{\text{equity }496.25}.
\]
The equity adjustment records the cumulative effect at the comparative opening. If the remaining lease year falls within the comparative year, the DTA of 27.5 and DTL of 26.25 reverse during that year, giving comparative deferred-tax expense 1.25. The opening adjustment alone therefore does not complete the transition. Neither ASBE 21's article 61(2) lease-adoption election nor Interpretation 16 Section III's no-comparative-adjustment rule for its separate share-payment topic may be substituted for Section I(b)'s tax transition.

These are gross recognition and measurement amounts. Separate recognition and noncurrent classification do not determine whether an actual balance sheet may offset DTA and DTL. That decision needs the applicable presentation provisions and entity/tax-authority facts. Actual tax law, tax bases and rates, recoverable profit, the enterprise's adoption and elections, and later amendments remain to be established before applying the example to an issuer. Full ASBE 22 credit, consolidated statements, group combinations and insurance revisions also retain specific source dependencies; similar amounts do not establish full IFRS equivalence.

\hypertarget{matched-case-outcomes-and-decision}{%
\section{Matched-case outcomes and decision}\label{method-standards-methods-matched-case-outcomes-and-decision}}

\begin{longtable}[]{@{}ll@{}}
\toprule
\begin{minipage}[b]{0.47\columnwidth}\raggedright
Input changed or withheld\strut
\end{minipage} & \begin{minipage}[b]{0.47\columnwidth}\raggedright
Required consequence\strut
\end{minipage}\tabularnewline
\midrule
\endhead
\begin{minipage}[t]{0.47\columnwidth}\raggedright
Framework changes with same lease facts\strut
\end{minipage} & \begin{minipage}[t]{0.47\columnwidth}\raggedright
Preserve different expense/ROU paths and source strength, not just rename accounts\strut
\end{minipage}\tabularnewline
\begin{minipage}[t]{0.47\columnwidth}\raggedright
Method source present but target entity's adoption unknown\strut
\end{minipage} & \begin{minipage}[t]{0.47\columnwidth}\raggedright
Unresolved applicability, no real posting authorization\strut
\end{minipage}\tabularnewline
\begin{minipage}[t]{0.47\columnwidth}\raggedright
ROU/liability initially equal\strut
\end{minipage} & \begin{minipage}[t]{0.47\columnwidth}\raggedright
Still evaluate gross tax differences and later divergent balances\strut
\end{minipage}\tabularnewline
\begin{minipage}[t]{0.47\columnwidth}\raggedright
Historical rule supports method but a later reporting year is requested\strut
\end{minipage} & \begin{minipage}[t]{0.47\columnwidth}\raggedright
Recheck operative amendments/transition; no automatic extension\strut
\end{minipage}\tabularnewline
\begin{minipage}[t]{0.47\columnwidth}\raggedright
Issuer policy discloses LIFO, regulator text unavailable\strut
\end{minipage} & \begin{minipage}[t]{0.47\columnwidth}\raggedright
Retain issuer assertion; do not upgrade it into universal US authority\strut
\end{minipage}\tabularnewline
\begin{minipage}[t]{0.47\columnwidth}\raggedright
Foreign basis label replaced by local label\strut
\end{minipage} & \begin{minipage}[t]{0.47\columnwidth}\raggedright
Reject as a standards conversion without recognition/measurement bridge\strut
\end{minipage}\tabularnewline
\bottomrule
\end{longtable}

These decisions establish a common procedure and specific proposed treatments. The recovered ASBE lease notice and tax interpretation support the stated application schedules and conditional example. Broader authority coverage remains incomplete: operative US sources, ASBE sector methods, actual tax law and presentation conditions, and the applicable amendment and entity-election evidence still require internal investigation.
